% Generated draft: requires mathematical and source review.
\documentclass[11pt]{article}
\usepackage[margin=25mm]{geometry}
\usepackage{fontspec}
\setmainfont{Latin Modern Roman}
\usepackage{amsmath,amssymb,mathtools,unicode-math}
\setmathfont{Latin Modern Math}
\usepackage{enumitem,needspace,xurl,hyperref}
\setlength{\parindent}{0pt}
\setlength{\parskip}{5pt}
\setlength{\emergencystretch}{3em}
\newcommand{\N}{\mathbb{N}}
\newcommand{\Z}{\mathbb{Z}}
\newcommand{\Q}{\mathbb{Q}}
\newcommand{\R}{\mathbb{R}}
\newcommand{\C}{\mathbb{C}}
\newcommand{\F}{\mathbb{F}}
\newcommand{\A}{\mathbb{A}}
\newcommand{\PP}{\mathbb{P}}
\newcommand{\E}{\mathbb{E}}
\newcommand{\eps}{\varepsilon}
\newcommand{\dd}{\,\mathrm{d}}
\begin{document}
\section*{UnsolvedMath ID: 30000315}
\textbf{Source problem code:} catalogue entry
\subsection{Definitions and mathematical statement}
Let $(M,\tau)$ be a finite von Neumann algebra with faithful normal finite trace, and let $A\subset M$ be a weakly dense $*$-subalgebra. The superscript in $M^A$ is essential. Its natural intended definition is the smallest $*$-subalgebra $B\subset M$ containing $A$ such that, for every abelian $*$-subalgebra $C\subset B$, the weak closure of $C$ in $M$ is also contained in $B$. The clean extract omits ``containing $A$'' in its short definition; that inclusion is needed to make the notation depend on $A$ and must be confirmed in the report. The first question asks for hypotheses ensuring
\[
 M^A=M.
\]
Weak density of $A$ alone is the setup, not a supplied proof of that equality.

Put $N=M\mathbin{\overline\otimes}M^{\mathrm{op}}$, with product trace, and let $U(N)$ be the algebra of closed densely defined operators affiliated with $N$, equipped with its measure topology. A basic convergence criterion is: $T_\alpha\to0$ in measure if for every $\varepsilon,\delta>0$, eventually there is a projection $e\in N$ with
\[
 (\tau\otimes\tau)(1-e)<\delta,\qquad\|T_\alpha e\|<\varepsilon.
\]
The domain measure topology uses the analogous definition in $M$.

Under the usual coarse bimodule convention, $a\cdot T\cdot b=(a\otimes b^{\mathrm{op}})T$. A complex-linear derivation $\Delta:M\to U(N)$ satisfies
\[
 \Delta(ab)=(a\otimes1)\Delta(b)+(1\otimes b^{\mathrm{op}})\Delta(a).
\]
The second question asks whether every such derivation is automatically continuous from bounded convergence in measure to convergence in measure: whenever a net $(x_\alpha)$ is uniformly bounded in operator norm and converges to $x$ in $\tau$-measure, must $\Delta(x_\alpha)\to\Delta(x)$ in the product-trace measure topology? No boundedness, closability or continuity of $\Delta$ is assumed in the question.

The precise bimodule, unitality and bounded-convergence conventions must agree with the report. In particular, $M^A$ is an algebra, not an ``associated map''. The literature assessment checked on 21 August 2026 describes an earlier underdefined extraction using $M_A$ and cannot by itself establish the status of the corrected algebra and derivation questions. The supplied verification restores those objects, while no later resolution is cited.
\subsection{Short English statement}
When does the abelian-weak-closure algebra generated from a weakly dense $A$ fill $M$? Separately, must every affiliated-operator-valued derivation carry uniformly bounded measure-convergent nets to measure-convergent nets? The restored superscript and bimodule conventions matter.
\subsection{Sources}
\begin{itemize}
\item \emph{C*-Algebren} (2006), Oberwolfach Reports, DOI 10.4171/owr/2005/41, pp. 2305--2374, \url{https://doi.org/10.4171/owr/2005/41}; proposed year supplied as 2005.
\item The sole source in the literature assessment checked 21 August 2026 is the original report, \url{https://doi.org/10.4171/owr/2005/41}. Its earlier ``map $M_A$'' wording is not treated as the definition of the corrected target.
\item Local UnsolvedMath registry, permanent ID $30000315$, exact source code \texttt{catalogue entry}, created 25 August 2026 and updated 16 September 2026. Supplied verification restores the superscripted algebra $M^A$, its intended definition and the derivation whose continuity is in question.
\end{itemize}
\end{document}
